\subsubsection{谱刚性与 Big-Witt 编译闭包：原子 Euler 数据的实现不变性、\texorpdfstring{$\gcd/\mathrm{lcm}$}{gcd/lcm} 乘法与 \texorpdfstring{$p$}{p}-典型伸缩}\label{subsec:operator-cyclic-traceclass-witt-rigidity}
小节 \ref{subsec:cyclic-traceclass-lift}、推论 \ref{cor:cyclic-euler-product-complex-section} 与定理 \ref{thm:cyclic-fredholm-realisation} 给出：对满足迹类可执行性（\(\ell^1\)）条件的原子 Euler 乘积，存在显式循环块直和算子 \(L_{\mathcal D}\) 使其 Fredholm \(\zeta\)-接口严格复原该乘积；而在正规实现子类中，这一模型在添零直和意义下已被完全分类。
本段记录四个互补的结构性后果：
(i) 任意迹类实现的\emph{非零谱多重集}由 determinant 精确强迫；
(ii) 循环块正规模型在相同 determinant 约束下极小化一切 Ky Fan 预算，并因而极小化所有酉不变范数；
(iii) 循环块在张量积下封闭，并以 \(\gcd/\mathrm{lcm}\) 的离散组合律实现 Big-Witt 乘法；
(iv) 在 \(p\)-幂长度族上可实现 Frobenius/Verschiebung 伸缩，给出 Dwork 型同余的谱实现模板。

\begin{proposition}[Fredholm 因式分解刚性：行列式约束强制精确非零谱]\label{prop:cyclic-euler-spectral-rigidity}\label{prop:fredholm-factorisation-rigidity}
\leanverified{paper\_cyclic\_euler\_spectral\_rigidity}
设
\[
\mathcal D=\{(\beta_j,n_j,m_j)\}_{j\in J},
\qquad
\mathcal D'=\{(\beta'_\ell,n'_\ell,m'_\ell)\}_{\ell\in L},
\]
为两组复权原子数据，满足
\[
\mathcal D,\mathcal D'
\subset (\CC\setminus\{0\})\times \NN_{\ge 1}\times \NN,
\qquad
\sum_{j\in J}m_jn_j|\beta_j|^{1/n_j}<\infty,
\qquad
\sum_{\ell\in L}m'_\ell n'_\ell|\beta'_\ell|^{1/n'_\ell}<\infty.
\]
\begin{enumerate}
  \item 若相应 Euler 乘积
  \[
  \prod_{j\in J}\bigl(1-\beta_j r^{n_j}\bigr)^{-m_j},
  \qquad
  \prod_{\ell\in L}\bigl(1-\beta'_\ell r^{n'_\ell}\bigr)^{-m'_\ell}
  \]
  在 \(r=0\) 的某个邻域内相等，则两组数据所对应的循环块正规模型具有完全相同的非零特征值多重集（按代数重数计）。特别地，
  \[
  \sum_{j\in J}m_jn_j=\sum_{\ell\in L}m'_\ell n'_\ell,
  \qquad
  \sum_{j\in J}m_jn_j|\beta_j|^{1/n_j}
  =
  \sum_{\ell\in L}m'_\ell n'_\ell|\beta'_\ell|^{1/n'_\ell}.
  \]
  \item 更一般地，若 \(T_1,T_2\in\mathcal{S}_1\) 满足
  \[
  \det(I-rT_1)^{-1}=\det(I-rT_2)^{-1}
  \]
  在 \(r=0\) 的某个邻域内成立，则 \(T_1,T_2\) 的非零谱多重集（按代数重数计）完全一致。因而
  \[
  \Tr(T_1^n)=\Tr(T_2^n),
  \qquad
  p_n(T_1)=p_n(T_2)
  \qquad (n\ge 1).
  \]
\end{enumerate}
\end{proposition}

\begin{proof}
对 \textup{(1)}，由定理 \ref{thm:cyclic-fredholm-realisation}，两组数据分别给出正规循环块模型 \(L_{\mathcal D}\) 与 \(L_{\mathcal D'}\)，并满足相应 Euler 乘积的 Fredholm 行列式 germ。取倒数后，
\[
\det(I-rL_{\mathcal D})=\det(I-rL_{\mathcal D'})
\]
在 \(r=0\) 邻域内成立，从而作为整函数恒等。由 \cite[Theorem~3.7]{Simon2005TraceIdeals}，\(\det(I-rT)\) 的零点恰为 \(T\) 的非零特征值倒数，且零点重数等于代数重数。故 \(L_{\mathcal D}\) 与 \(L_{\mathcal D'}\) 具有完全相同的非零谱多重集；两条显示公式随之立即得到。

对 \textup{(2)}，把上述论证直接应用于 \(T_1,T_2\) 即得其非零谱多重集完全相同。再比较
\[
\log\det(I-rT_i)^{-1}
=\sum_{n\ge 1}\frac{\Tr(T_i^n)}{n}r^n
\qquad (i=1,2)
\]
的系数，即得迹序列相同；再经 Möbius 反演便有 \(p_n(T_1)=p_n(T_2)\)。证毕。
\end{proof}

\begin{corollary}[行列式约束下的精确非零谱]\label{thm:general-support-optimality}
设 \(T\in\mathcal{S}_1(H)\) 满足
\[
\det(I-rT)^{-1}
=\prod_{j\in J}\bigl(1-\beta_j r^{n_j}\bigr)^{-m_j}
\]
在 \(r=0\) 的某个邻域内成立，其中
\[
\mathcal D=\{(\beta_j,n_j,m_j)\}_{j\in J}
\subset (\CC\setminus\{0\})\times \NN_{\ge 1}\times \NN,
\qquad
\sum_{j\in J}m_jn_j|\beta_j|^{1/n_j}<\infty.
\]
任选 \(\alpha_j\in\CC\) 使 \(\alpha_j^{n_j}=\beta_j\)。则 \(T\) 的非零谱多重集（按代数重数计）恰为
\[
\biguplus_{j\in J}\bigl\{\alpha_j\omega:\omega^{n_j}=1\bigr\}^{m_j}.
\]
特别地，其非零谱总代数重数为
\[
\sum_{j\in J}m_jn_j.
\]
\end{corollary}

\begin{proof}
设 \(L_{\mathcal D}\) 为定理 \ref{thm:cyclic-fredholm-realisation} 所给的循环块正规模型。由命题 \ref{prop:fredholm-factorisation-rigidity}\textup{(2)}，\(T\) 与 \(L_{\mathcal D}\) 具有完全相同的非零谱多重集；而 \(L_{\mathcal D}\) 的非零谱正是各个循环块的单位根轨道并置。证毕。
\end{proof}

\begin{corollary}[Schatten--\texorpdfstring{$p$}{p} 范数的原子闭式与实现不变性（正规实现）]\label{cor:cyclic-euler-schatten-invariants}
设
\[
\mathcal D=\{(\beta_j,n_j,m_j)\}_{j\in J}
\subset (\CC\setminus\{0\})\times \NN_{\ge 1}\times \NN
\]
满足
\[
\sum_{j\in J}m_jn_j|\beta_j|^{1/n_j}<\infty,
\]
并设 \(L\) 为对应 Euler 乘积的一个\emph{正规}迹类实现（例如 \(L=L_{\mathcal D}\) 的循环块直和实现）。则对任意 \(p>0\)，其 Schatten--\(p\)（准）范数满足无条件闭式
\[
\|L\|_{S_p}^{p}
=\sum_{j\in J} m_j n_j\,|\beta_j|^{p/n_j}.
\]
因此 \(\{\|L\|_{S_p}\}_{p>0}\) 构成对原子 Euler 数据的可计算谱指纹族，即使不同原子贡献的单位根轨道发生重合也不例外。
\end{corollary}

\begin{proof}
由定理 \ref{thm:cyclic-fredholm-realisation}，任意正规实现都与某个 \(0\oplus L_{\mathcal D}\) 酉等价；等价地，由推论 \ref{thm:general-support-optimality}，其非零特征值多重集正是
\(
\biguplus_j\{\alpha_j\omega:\omega^{n_j}=1\}^{m_j}
\)。
对正规算子，奇异值恰为特征值模长（按重数计），而每个 \(j\) 恰贡献 \(m_jn_j\) 个模长 \(|\beta_j|^{1/n_j}\) 的非零特征值。故
\[
\|L\|_{S_p}^{p}
=\sum_{j\in J} m_j n_j\,|\beta_j|^{p/n_j}.
\]
证毕。
\end{proof}

\begin{theorem}[行列式约束下的 Ky Fan 优势与酉不变范数极小性]\label{thm:kyfan-dominance}
设 \(T\in\mathcal{S}_1(H)\) 满足
\[
\det(I-rT)^{-1}
=\prod_{j\in J}\bigl(1-\beta_j r^{n_j}\bigr)^{-m_j}
\]
在 \(r=0\) 的某个邻域内成立，其中
\[
\mathcal D=\{(\beta_j,n_j,m_j)\}_{j\in J}
\subset (\CC\setminus\{0\})\times \NN_{\ge 1}\times \NN
\]
满足
\[
\sum_{j\in J}m_jn_j|\beta_j|^{1/n_j}<\infty.
\]
令 \(L_{\mathcal D}\) 为定理 \ref{thm:cyclic-fredholm-realisation} 的循环块正规模型，并记奇异值按非增序排为
\[
s_1(T)\ge s_2(T)\ge \cdots,
\qquad
s_1(L_{\mathcal D})\ge s_2(L_{\mathcal D})\ge \cdots.
\]
则对任意 \(N\ge 1\)，都有
\[
\sum_{k=1}^{N}s_k(L_{\mathcal D})
\le
\sum_{k=1}^{N}s_k(T).
\]
因此，在所有具有同一 Fredholm determinant 的迹类实现中，\(L_{\mathcal D}\) 同时最小化一切 Ky Fan 范数；进而由 Ky Fan dominance 原理，凡在 \(L_{\mathcal D}\) 上有限的酉不变范数 \(\nu\) 都满足
\[
\nu(L_{\mathcal D})\le \nu(T).
\]
\end{theorem}

\begin{proof}
由推论 \ref{thm:general-support-optimality}，\(T\) 与 \(L_{\mathcal D}\) 具有完全相同的非零特征值多重集。将这些非零特征值按模长非增序记为
\[
\lambda_1,\lambda_2,\dots.
\]
由于 \(L_{\mathcal D}\) 正规，其奇异值恰为这些特征值的模长：
\[
s_k(L_{\mathcal D})=|\lambda_k|
\qquad (k\ge 1).
\]
经典的 Ky Fan 特征值--奇异值不等式给出
\[
\sum_{k=1}^{N}|\lambda_k|
\le
\sum_{k=1}^{N}s_k(T)
\qquad (N\ge 1),
\]
与上一式合并便得所述 Ky Fan 优势。最后一段关于任意酉不变范数的结论，是 Ky Fan dominance 原理对上述偏和不等式的直接推论。证毕。
\end{proof}

\begin{theorem}[迹范数极小性与正规刻画]\label{thm:trace-norm-minimality}
设 \(T\in\mathcal{S}_1(H)\) 满足
\[
\det(I-rT)^{-1}
=\prod_{j\in J}\bigl(1-\beta_j r^{n_j}\bigr)^{-m_j}
\]
在 \(r=0\) 的某个邻域内成立，其中
\[
\mathcal D=\{(\beta_j,n_j,m_j)\}_{j\in J}
\subset (\CC\setminus\{0\})\times \NN_{\ge 1}\times \NN
\]
满足
\[
\sum_{j\in J}m_jn_j|\beta_j|^{1/n_j}<\infty.
\]
则
\[
\|T\|_1
\ge
\sum_{j\in J}m_jn_j|\beta_j|^{1/n_j}.
\]
等号成立当且仅当 \(T\) 正规。等价地，等号成立当且仅当存在 Hilbert 空间 \(H_0\) 与酉算子 \(U:H\to H_0\oplus H_{\mathcal D}\)，使
\[
UTU^{-1}=0_{H_0}\oplus L_{\mathcal D},
\]
其中 \(L_{\mathcal D}\) 是定理 \ref{thm:cyclic-fredholm-realisation} 的循环块正规模型。
\end{theorem}

\begin{proof}
令 \(L_{\mathcal D}\) 为定理 \ref{thm:cyclic-fredholm-realisation} 所给的循环块正规模型。由定理 \ref{thm:kyfan-dominance} 对 \(N\to\infty\) 取极限，得到
\[
\|L_{\mathcal D}\|_1\le \|T\|_1.
\]
再由定理 \ref{thm:cyclic-fredholm-realisation} 对正规模型的迹范数公式，
\[
\|L_{\mathcal D}\|_1
=\sum_{j\in J}m_jn_j|\beta_j|^{1/n_j},
\]
于是得出所述下界。

若 \(T\) 正规，则由定理 \ref{thm:cyclic-fredholm-realisation}，\(T\) 与某个 \(0_{H_0}\oplus L_{\mathcal D}\) 酉等价，故显然取到等号。反过来设 \(\|T\|_1=\|L_{\mathcal D}\|_1\)。将 \(T\) 的非零特征值按代数重数重复并按模长非增序记为
\[
\lambda_1(T),\lambda_2(T),\dots.
\]
由推论 \ref{thm:general-support-optimality}，
\[
\sum_k|\lambda_k(T)|
=\sum_{j\in J}m_jn_j|\beta_j|^{1/n_j}
=\|T\|_1.
\]

取一组将紧算子 \(T\) 上三角化的正交基，使 \(T\) 的对角元恰为 \(\lambda_k(T)\)（按代数重数重复）。定义对角酉算子 \(D\) 的对角元为
\[
d_k=
\begin{cases}
\overline{\lambda_k(T)}/|\lambda_k(T)|,
& \lambda_k(T)\neq 0,\\
1,& \lambda_k(T)=0.
\end{cases}
\]
则 \(DT\) 仍为上三角，且
\[
\Tr(DT)=\sum_k|\lambda_k(T)|=\|T\|_1.
\]
把 \(T\) 写成极分解 \(T=V|T|\)。由于 \(\|DV\|\le 1\)，有
\[
0
=\Tr|T|-\Re\Tr(DT)
=\Tr\!\bigl(|T|^{1/2}(I-\Re(DV))|T|^{1/2}\bigr).
\]
括号中的算子是正的，故必为零。于是对任意
\(
h\in\overline{\operatorname{ran}(|T|)}
\)
都有 \(\Re\langle DVh,h\rangle=\|h\|^2\)，从而
\[
\|h-DVh\|^2
=\|h\|^2+\|DVh\|^2-2\Re\langle DVh,h\rangle
\le 2\|h\|^2-2\|h\|^2=0.
\]
故 \(DVh=h\)，于是
\[
DT=DV|T|=|T|
\]
为正算子。正的上三角算子只能是对角算子，因此 \(DT\) 为对角算子，进而 \(T\) 也在该基下对角化，故 \(T\) 正规。最后一条等价表述再由定理 \ref{thm:cyclic-fredholm-realisation} 立即得到。证毕。
\end{proof}

\begin{proposition}[循环块张量积的分解定理：\texorpdfstring{$\gcd/\mathrm{lcm}$}{gcd/lcm} 控制的 Witt 乘法实现]\label{prop:cyclic-block-tensor-gcd-lcm}
令 \(C(n,\alpha)=\alpha\Pi_n\) 为定义 \ref{def:cyclic-block} 的循环块。
对任意 \(n,n'\ge 1\)，置
\[
\ell:=\operatorname{lcm}(n,n'),\qquad g:=\gcd(n,n').
\]
则存在显式置换酉 \(U\) 使得在 \(\CC^{n}\otimes\CC^{n'}\) 上成立酉等价
\[
\boxed{\;
C(n,\alpha)\otimes C(n',\alpha')\ \simeq\
\bigoplus_{i=1}^{g} C(\ell,\alpha\alpha')\; }.
\]
\leanverified{paper\_cyclic\_block\_tensor\_gcd\_lcm\_seeds}
\end{proposition}

\begin{proof}[证明要点]
算子 \(\Pi_n\otimes\Pi_{n'}\) 是有限集合 \(\ZZ/n\ZZ\times\ZZ/n'\ZZ\) 上“同步平移 \((a,b)\mapsto(a+1,b+1)\)”的置换矩阵。
该置换的轨道分解为 \(g=\gcd(n,n')\) 条不交轨道，每条轨道长度为 \(\ell=\operatorname{lcm}(n,n')\)；
因此存在置换基变换 \(U\) 使
\(
U(\Pi_n\otimes\Pi_{n'})U^{-1}=\bigoplus_{i=1}^{g}\Pi_{\ell}
\)。
再乘以标量 \(\alpha\alpha'\) 得到所示分解。证毕。
\end{proof}

\begin{corollary}[张量积因子的行列式闭式]\label{cor:cyclic-block-tensor-determinant}
在命题 \ref{prop:cyclic-block-tensor-gcd-lcm} 的设定下，对任意 \(r\in\CC\) 有
\[
\det\!\bigl(I-r\,(C(n,\alpha)\otimes C(n',\alpha'))\bigr)
=\bigl(1-(\alpha\alpha' r)^{\ell}\bigr)^{g}.
\]
若进一步设 \(\beta=\alpha^n\)、\(\beta'=(\alpha')^{n'}\)，则等价地
\[
\det\!\bigl(I-r\,(C(n,\alpha)\otimes C(n',\alpha'))\bigr)
=\bigl(1-\beta^{\ell/n}(\beta')^{\ell/n'}\,r^{\ell}\bigr)^{g}.
\]
\end{corollary}

\begin{proof}
由酉等价与行列式对直和的乘法性、以及命题 \ref{prop:cycle-permutation-determinant} 的恒等式
\(
\det(I-t\Pi_{\ell})=1-t^{\ell}
\)
立即得到第一式；第二式由 \((\alpha\alpha')^{\ell}=\alpha^{\ell}(\alpha')^{\ell}=\beta^{\ell/n}(\beta')^{\ell/n'}\)（指数为整数）化简得到。证毕。
\end{proof}

\begin{proposition}[\texorpdfstring{$p$}{p}-典型 Frobenius/Verschiebung 的算子实现：循环长度伸缩与 ghost 兼容]\label{prop:cyclic-p-typical-frobenius-verschiebung}
\leanverified{paper\_cyclic\_p\_typical\_frobenius\_verschiebung}
固定素数 \(p\)。
在 \(p\)-幂循环块族 \(\{C(p^k,\alpha)\}_{k\ge 0}\) 上定义两种操作：
\[
\mathsf F_p\bigl(C(p^k,\alpha)\bigr):=
\begin{cases}
\bigoplus_{i=1}^{p} C(p^{k-1},\alpha^{p}),& k\ge 1,\\
C(1,\alpha^{p}),& k=0,
\end{cases}
\qquad
\mathsf V_p\bigl(C(p^k,\alpha)\bigr):=C(p^{k+1},\widetilde\alpha)\ \ (\widetilde\alpha^{p}=\alpha),
\]
并按直和逐块延拓到有限直和算子。
则对任意由 \(p\)-幂循环块有限直和得到的算子 \(L\)，对一切 \(n\ge 1\) 都有 ghost 兼容性
\[
\boxed{\ \Tr\!\bigl(\mathsf F_p(L)^n\bigr)=\Tr\!\bigl(L^{pn}\bigr)\ }.
\]
并且在酉等价类与直和加法的意义下，
\[
\boxed{\ \mathsf F_p\!\bigl(\mathsf V_p(L)\bigr)\ \simeq\ L^{\oplus p}\ }.
\]
因此 \(\mathsf F_p\circ \mathsf V_p\) 在该 \(p\)-典型分量上实现“乘以 \(p\)”的 Verschiebung--Frobenius 关系。
\end{proposition}

\begin{proof}[证明要点]
对单块 \(C(p^k,\alpha)\)，由 \(\Tr(\Pi_{p^k}^m)=p^k\mathbf{1}_{p^k\mid m}\) 得
\[
\Tr\!\bigl(C(p^k,\alpha)^{m}\bigr)=p^k\,\alpha^{m}\,\mathbf{1}_{p^k\mid m}.
\]
据此直接计算可验证：对任意 \(n\ge 1\)，
\(\Tr(\mathsf F_p(C(p^k,\alpha))^n)=\Tr(C(p^k,\alpha)^{pn})\)；
对直和按迹的可加性成立。
第二式由 \((\widetilde\alpha)^p=\alpha\) 与定义立得：
\(\mathsf F_p(C(p^{k+1},\widetilde\alpha))=\bigoplus_{i=1}^{p}C(p^k,\alpha)\)。
证毕。
\end{proof}

\begin{theorem}[Witt--算子函子：原子 Euler 数据的 Big-Witt 运算可编译为迹类算子组合]\label{thm:atomic-witt-into-tc1}
令 \(\mathrm{TC}_1\) 表示迹类算子在“酉等价类”下的交换半环：加法为直和，乘法为张量积。
令 \(\mathcal{W}_{\mathrm{atom}}\) 为由原子数据 \((\beta,n,m)\) 生成的 Big-Witt 型半环，其中加法对应 Euler 乘积相乘（原子多重集并置），乘法由 \(\gcd/\operatorname{lcm}\) 规则给出。
则映射
\[
(\beta,n,m)\ \longmapsto\ \bigoplus_{i=1}^{m} C\!\left(n,\alpha\right),
\qquad \alpha^{n}=\beta,
\]
在加法上对应直和，并在乘法上由命题 \ref{prop:cyclic-block-tensor-gcd-lcm} 保证与 \(\gcd/\operatorname{lcm}\) 规则严格一致；
因此其延拓为一个构造性的半环同态
\[
\boxed{\ \mathcal{W}_{\mathrm{atom}}\ \longrightarrow\ \mathrm{TC}_1\ }.
\]
从而 Big-Witt 的抽象运算可被编译为迹类算子的可审计组合运算（直和/张量积），并与推论 \ref{cor:cyclic-euler-product-complex-section} 的 Fredholm \(\zeta\)-接口严格相容。
\end{theorem}

\begin{theorem}[Fredholm \texorpdfstring{$\zeta$}{zeta} 界面的等同性判定不可判定]\label{thm:operator-fredholm-zeta-equality-undecidable}
存在一类满足迹类可执行性条件的原子数据流 \(\mathcal D\)（以可计算方式生成），使得如下判定问题不可判定：
给定两份可计算原子数据流 \(\mathcal D_0,\mathcal D_1\)，判定其 Euler 乘积
\[
\zeta_{\mathcal D_0}(r)=\prod_{j}(1-\beta^{(0)}_j r^{n^{(0)}_j})^{-m^{(0)}_j},
\qquad
\zeta_{\mathcal D_1}(r)=\prod_{j}(1-\beta^{(1)}_j r^{n^{(1)}_j})^{-m^{(1)}_j}
\]
在某个圆盘 \(|r|<r_0\) 上作为解析函数是否恒等相等。
等价地，借助循环块直和实现（推论 \ref{cor:cyclic-euler-product-complex-section}），判定
\(
\det(I-rL_{\mathcal D_0})^{-1}\equiv \det(I-rL_{\mathcal D_1})^{-1}
\)
是否成立亦不可判定。
\leanverified{fredholm\_equality\_decidable\_finite}
\end{theorem}

\begin{proof}
对任意图灵机 \(M\) 与输入 \(x\)，构造一份可计算原子数据流 \(\mathcal D_{M,x}\)：
按步数 \(t=1,2,\dots\) 枚举并模拟 \(M(x)\) 至多 \(t\) 步；若在第 \(t\) 步首次停机，则输出单个原子因子 \((\beta,n,m)=(2^{-t},t,1)\)，并在其后输出空数据；若未停机则始终输出空数据。
该流显然满足迹类可执行性（至多输出一个因子），且
\[
\zeta_{\mathcal D_{M,x}}(r)=
\begin{cases}
(1-2^{-t}r^{t})^{-1},& M(x)\ \text{于第 }t\text{ 步停机},\\
1,& M(x)\ \text{不停机}.
\end{cases}
\]
因此判定 \(\zeta_{\mathcal D_{M,x}}(r)\equiv 1\) 等价于判定 \(M(x)\) 是否停机，从而不可判定。
将 \(\mathcal D_0=\mathcal D_{M,x}\)、\(\mathcal D_1=\varnothing\)（空数据，对应 \(\zeta\equiv 1\)）代入即得结论。
Fredholm 行列式表述由推论 \ref{cor:cyclic-euler-product-complex-section} 把 \(\zeta_{\mathcal D}\) 实现为 \(\det(I-rL_{\mathcal D})^{-1}\) 立即推出。证毕。
\end{proof}

\begin{remark}[卷积--Fourier 对角化与 Witt--ghost 的统一解释]\label{rem:operator-fourier-witt-unification}
注 \ref{rem:twist-commuting-cube} 已表明：折叠多重度的角色扭曲 \(\mathrm{Tw}_\chi\) 同时对齐 Fourier 对角化与 Witt/Euler 交换图。
定理 \ref{thm:atomic-witt-into-tc1} 进一步给出一个可执行的“算子化提升”观点：在原子闭包子类上，可将 ghost 坐标（迹序列）视为某个迹类算子的 \(\Tr(L^n)\)，
从而把“卷积--Fourier 的乘积分解”与“Witt--ghost 的坐标变换”统一为同一对象在不同函子下的读出。
该统一口径为把折叠计数流接入 Fredholm \(\zeta\) 接口提供了最小的可审计编译模板。
\end{remark}

\endinput
